\documentclass[10pt,a4paper]{amsart}
\usepackage[margin=19mm]{geometry}
\usepackage{amsmath,amssymb,mathtools}
\usepackage[scr=rsfs,cal=euler]{mathalfa}
\usepackage{xfrac}
\usepackage[unicode,hidelinks,bookmarks=false]{hyperref}
\newcommand{\ie}{i.e.\ }
\newcommand{\cf}{cf.\ }
\newcommand{\resp}{respectively\ }
\newcommand{\Subsection}{\subsection}
\providecommand{\nobreakdash}{\nobreak\hbox{-}}
\setlength{\emergencystretch}{2em}
\multlinegap0pt
\usepackage{amssymb}
\newtheorem{lemma}{Lemma}
\newtheorem{proposition}{Proposition}
\title{An adjunction inequality for Real embedded surfaces}
\author{David Baraglia}
\date{Source: arXiv:2507.05667v2 (CC BY 4.0)}
\begin{document}
\maketitle

\noindent\emph{Source and licence.} An adjunction inequality for Real embedded surfaces. Version arXiv:2507.05667v2 (CC BY 4.0), distributed by arXiv under the Creative Commons Attribution 4.0 International licence, http://creativecommons.org/licenses/by/4.0/, as recorded in the arXiv licence metadata for this exact version and shown on the abstract page https://arxiv.org/abs/2507.05667v2. Full licence notice: \texttt{licenses/arxiv-2507.05667-CC-BY-4.0.txt}.\par\smallskip

\noindent\textit{Editorial scope.} Counted unit: the article's proposition prop:deg2 (source0.tex lines 263--272). The article's complete original proof of it is delivered with it (source0.tex lines 273--287, one \texttt{proof} environment of 15 lines). No mathematics is written, repaired or re-derived: the statement and the proof are the article's own lines, in the article's own notation, and no retained line is substituted or re-flowed.  Retained uncounted as the source-local context the proof needs: lines 254--256.  Nothing was omitted from any retained span, so the package carries no omission record.  No retained line contains a citation, so the wrapper carries no bibliography and no bibliography entry.  No retained line contains a figure environment, an image-inclusion command or drawing code, so no image is delivered and the package declares no asset dependency.  Editorial formatting only, in addition: an AMS \texttt{amsart} preamble that restates the article's own declarations and macros used by the retained lines, namely the environments \texttt{lemma}, \texttt{proposition} on separate counters, together with the standard packages those lines need.  The retained environments print in this document as Lemma 1, Proposition 1 in order of appearance, whereas the article prints the same items with the numbering of its own full document.  The complete native source of the article as distributed in the arXiv e-print for this exact version is bundled unchanged as \texttt{sources/181195/source0.tex}.
\begin{lemma}\label{lem:rdivgen}
Every Real line bundle on $\Sigma$ is isomorphic to a Real line bundle of the form $\mathcal{O}(D)$ for a Real divisor $D$.
\end{lemma}

\begin{proposition}\label{prop:deg2}
The homomorphism $\phi : H^2_{\mathbb{Z}_2}(\Sigma ; \mathbb{Z}_-) \to \mathbb{Z} \oplus \mathbb{Z}_2^n$ given by
\[
\phi(E) = ( deg(E) , u_1(E) , \dots , u_n(E))
\]
is injective with image equal to 
\[
A = \{ (d , u_1 , \dots , u_n) \in \mathbb{Z} \oplus \mathbb{Z}_2^n \; | \; d = u_1 + \cdots + u_n \; ({\rm mod} \; 2) \}.
\]
\end{proposition}

\begin{proof}
We first show that the image of $\phi$ is $A$. Let $x_i \in C_i$. Then $\phi( \mathcal{O}(x_i) ) = (1 , e_i )$, where $e_i$ denotes the $i$-th standard basis vector of $\mathbb{Z}_2^n$. Let $x \in \Sigma$ be a non-fixed point. Then $\phi( \mathcal{O}(x+\sigma(x)) ) = (2 , 0 , \dots , 0)$. Clearly the elements $(1 , e_i)$, $i = 1, \dots , n$ and $(2 , 0 , \dots , 0)$ belong to $A$ and in fact generate $A$. So the image of $\phi$ contains $A$. Since any Real divisor can be written as a linear combination of divisors of the form $x_i$, $x_i \in C_i$ or $x + \sigma(x)$, $\sigma(x) \neq x$, we see that $\phi( \mathcal{O}(D) ) \in A$ for any Real divisor $D$. Then by Lemma \ref{lem:rdivgen}, it follows that the image of $\phi$ is $A$.

Now we show $\phi$ is injective. If $n = 0$, then $\sigma$ acts freely. Let $M = \Sigma/\sigma$. Then $\Sigma \to M$ is the orientation double cover of $M$ and the equivariant local system $\mathbb{Z}_-$ descends via $\sigma$ to the orientation local system $\mathbb{Z}_{orn}$ on $M$. Hence
\[
H^2_{\mathbb{Z}_2}(\Sigma ; \mathbb{Z}_-) \cong H^2(M ; \mathbb{Z}_{orn}) \cong H_0(M ; \mathbb{Z}) \cong \mathbb{Z}.
\]
If $n=0$, then $A \cong \mathbb{Z}$. Hence $\phi$ can be indentified with a surjective homomorphism $\phi : \mathbb{Z} \to \mathbb{Z}$. This implies that $\phi$ is injective.

If $n > 0$, then the Borel spectral sequence yield a short exact sequence
\[
0 \to \mathbb{Z}_2^c \to H^2_{\mathbb{Z}_2}(\Sigma ; \mathbb{Z}_-) \to \mathbb{Z} \to 0.
\]
This sequence splits, so $H^2_{\mathbb{Z}_2}(\Sigma ; \mathbb{Z}_-) \cong \mathbb{Z} \oplus \mathbb{Z}_2^c$. Using localisation in equivariant cohomology one finds that $c = n-1$, so $H^2_{\mathbb{Z}_2}(\Sigma ; \mathbb{Z}_-) \cong \mathbb{Z} \oplus \mathbb{Z}_2^{n-1}$. Similarly, one sees that $A \cong \mathbb{Z} \oplus \mathbb{Z}_2^{n-1}$, so $\phi$ can be identified with a surjective homomorphism $\phi : \mathbb{Z} \oplus \mathbb{Z}_2^{n-1} \to \mathbb{Z} \oplus \mathbb{Z}_2^{n-1}$. Such a homomorphism is easily seen to be an isomorphism.
\end{proof}

\end{document}
